\documentclass[10pt,a4paper]{amsart}
\usepackage[margin=19mm]{geometry}
\usepackage{amsmath,amssymb,mathtools}
\usepackage[scr=rsfs,cal=euler]{mathalfa}
\usepackage{xfrac}
\usepackage[unicode,hidelinks,bookmarks=false]{hyperref}
\newcommand{\ie}{i.e.\ }
\newcommand{\cf}{cf.\ }
\newcommand{\resp}{respectively\ }
\newcommand{\Subsection}{\subsection}
\providecommand{\nobreakdash}{\nobreak\hbox{-}}
\setlength{\emergencystretch}{2em}
\multlinegap0pt
\usepackage{amssymb}
\usepackage{dsfont}
\usepackage[normalem]{ulem}
\usepackage{tikz}
\usepackage{mathtools}
\usepackage{enumerate}
\usepackage{xcolor}
\newtheorem{lem}{Lemma}
\newtheorem{prop}{Proposition}
\newtheorem{thm}{Theorem}
\title{Holomorphic Anomaly Equations For $[\mathbb{C}^5/\mathbb{Z}_5]$}
\author{Deniz Genlik, Hsian-Hua Tseng}
\date{Source: arXiv:2211.15878v3 (CC BY 4.0)}
\begin{document}
\maketitle

\noindent\emph{Source and licence.} Holomorphic Anomaly Equations For $[\mathbb{C}^5/\mathbb{Z}_5]$. Version arXiv:2211.15878v3 (CC BY 4.0), distributed by arXiv under the Creative Commons Attribution 4.0 International licence, http://creativecommons.org/licenses/by/4.0/, as recorded in the arXiv licence metadata for this exact version and shown on the abstract page https://arxiv.org/abs/2211.15878v3. Full licence notice: \texttt{licenses/arxiv-2211.15878-CC-BY-4.0.txt}.\par\smallskip

\noindent\textit{Editorial scope.} Counted unit: the article's thm thm:HAE\_wrt\_A2\_partial (source0.tex lines 632--639). The article's complete original proof of it is delivered with it (source0.tex lines 640--718, one \texttt{proof} environment of 79 lines). No mathematics is written, repaired or re-derived: the statement and the proof are the article's own lines, in the article's own notation, and no retained line is substituted or re-flowed.  Retained uncounted as the source-local context the proof needs: lines 210--212; lines 380--382; lines 526--548; lines 551--561; lines 588--594.  Nothing was omitted from any retained span, so the package carries no omission record.  No retained line contains a citation, so the wrapper carries no bibliography and no bibliography entry.  No retained line contains a figure environment, an image-inclusion command or drawing code, so no image is delivered and the package declares no asset dependency.  Editorial formatting only, in addition: an AMS \texttt{amsart} preamble that restates the article's own declarations and macros used by the retained lines, namely the environments \texttt{lem}, \texttt{prop}, \texttt{thm} on separate counters, together with the standard packages those lines need.  The retained environments print in this document as Lemma 1, Proposition 1, Theorem 1 in order of appearance, whereas the article prints the same items with the numbering of its own full document.  The complete native source of the article as distributed in the arXiv e-print for this exact version is bundled unchanged as \texttt{sources/134149/source0.tex}.
\begin{equation}\label{eq:C12C22C3L5}
C_1^2C_2^2C_3=L^5.
\end{equation}

\begin{lem}\label{lem:P0jk_is_in_CLplusminus}
We have $\widetilde{P}_{0,j}^k\in\mathbb{C}[L^{\pm 1}]$ for all $0\leq{j}\leq{4}$ and $k\geq{0}$.
\end{lem}

\begin{prop}\label{prop:contributions}
The contribution associated to a decorated stable graph $\Gamma\in\mathrm{G}_{g,n}^{\text{Dec}}(5)$ is
\begin{equation*}
\operatorname{Cont}_{\Gamma}\left(\phi_{c_{1}}, \ldots, \phi_{c_{n}}\right)=\frac{1}{|\mathrm{Aut}(\Gamma^{\mathrm{St}})|} \sum_{\mathrm{A} \in \mathbb{Z}_{\geq 0}^{\mathrm{F}(\Gamma)}} \prod_{\mathfrak{v} \in \mathrm{V}_{\Gamma}} \mathrm{Cont}_{\Gamma}^{\mathrm{A}}(\mathfrak{v}) \prod_{\mathfrak{e}\in \mathrm{E}_{\Gamma}} \mathrm{Cont}_{\Gamma}^{\mathrm{A}}(\mathfrak{e}) \prod_{\mathfrak{l} \in \mathrm{L}_{\Gamma}} \mathrm{Cont}_{\Gamma}^{\mathrm{A}}(\mathfrak{l})
\end{equation*}
where $\mathrm{F}(\Gamma)=\left\vert\mathrm{H}_{\Gamma}\right\vert$. Here, $\mathrm{Cont}_{\Gamma}^{\mathrm{A}}(\mathfrak{v})$, $\mathrm{Cont}_{\Gamma}^{\mathrm{A}}(\mathfrak{e})$, and $\mathrm{Cont}_{\Gamma}^{\mathrm{A}}(\mathfrak{l})$ are the {\em vertex}, {\em edge} and {\em leg} contributions with flag $\mathrm{A}-$values\footnote{Notation: The values ${b_{\mathfrak{v}1}},\ldots,{b_{\mathfrak{v}\mathrm{h}(\mathfrak{v})}}$ and $b_{\mathfrak{e}1},b_{\mathfrak{e}2}$ are the entries of $(a_1,\ldots,a_m,b_{m+1},\ldots,b_{\left\vert\mathrm{H}_{\Gamma}\right\vert})$ corresponding to $\mathrm{Cont}_{\Gamma}^{\mathrm{A}}(\mathfrak{v})$ and $\mathrm{Cont}_{\Gamma}^{\mathrm{A}}(\mathfrak{e})$ respectively.} $(a_1,\ldots,a_m,b_{m+1},\ldots,b_{\left\vert\mathrm{H}_{\Gamma}\right\vert})$ respectively, and they are given by
\begin{equation*}
\begin{aligned}
    \mathrm{Cont}_{\Gamma}^{\mathrm{A}}(\mathfrak{v})
    =&\sum_{k \geq 0} \frac{g({e}_{\mathrm{p}(\mathfrak{v})},{e}_{\mathrm{p}(\mathfrak{v})})^{-\frac{2\mathrm{g}(\mathfrak{v})-2+\mathrm{n}(\mathfrak{v})+k}{2}}}{k !}\\
    &\times\int_{\overline{M}_{\mathrm{g}(\mathfrak{v}),\mathrm{n}(\mathfrak{v})+k}}\psi_1^{a_{\mathfrak{v}1}}\cdots\psi_{\mathrm{l}(v)}^{a_{\mathfrak{v}\mathrm{l}(\mathfrak{v})}}\psi_{\mathrm{l}(\mathfrak{v})+1}^{b_{\mathfrak{v}1}}\cdots\psi_{\mathrm{n}(\mathfrak{v})}^{b_{\mathfrak{v}\mathrm{h}(\mathfrak{v})}}t_{\mathrm{p}(\mathfrak{v})}(\psi_{\mathrm{n}(\mathfrak{v})+1})\cdots t_{\mathrm{p}(\mathfrak{v})}(\psi_{\mathrm{n}(\mathfrak{v})+k}),\\
    \mathrm{Cont}_{\Gamma}^{\mathrm{A}}(\mathfrak{e})
    =&\frac{(-1)^{b_{\mathfrak{e}1}+b_{\mathfrak{e}2}}}{5} \sum_{m=0}^{b_{\mathfrak{e}2}}(-1)^{m} \sum_{r=0}^{4}\frac{\widetilde{P}_{\mathrm{Inv}(r),\mathrm{p}(\mathfrak{v}_1)}^{b_{\mathfrak{e}1}+m+1}\widetilde{P}_{r,\mathrm{p}(\mathfrak{v}_2)}^{b_{\mathfrak{e}2}-m}}{\zeta^{(b_{\mathfrak{e}1}+m+1+\mathrm{Inv}(r))\mathrm{p}(\mathfrak{v}_1)}\zeta^{(b_{\mathfrak{e}2}-m+r)\mathrm{p}(\mathfrak{v}_2)}},\\
    \mathrm{Cont}_{\Gamma}^{\mathrm{A}}(\mathfrak{l})
    =&\frac{(-1)^{a_{\ell(\mathfrak{l})}}}{5}\frac{K_{\mathrm{Inv}(c_{\ell(\mathfrak{l})})}}{L^{\mathrm{Inv}(c_{\ell(\mathfrak{l})})}}
    \frac{\widetilde{P}_{\mathrm{Inv}(c_{\ell(\mathfrak{l})}),\mathrm{p}(\nu(\mathfrak{l}))}^{{a_{\ell(\mathfrak{l})}}}}{     \zeta^{({a_{\ell(\mathfrak{l})}}+{\mathrm{Inv}(c_{\ell(\mathfrak{l})})})\mathrm{p}(\nu(\mathfrak{l}))}},
\end{aligned}
\end{equation*}
where
\begin{equation*}
t_{\mathrm{p}(\mathfrak{v})}(z)=\sum_{j\geq{2}}\mathrm{T}_{\mathrm{p}(\mathfrak{v})j}z^j\quad\text{with}\quad \mathrm{T}_{\mathrm{p}(\mathfrak{v})j}=\frac{(-1)^j}{n}\widetilde{P}_{0,\mathrm{p}(\mathfrak{v})}^{j-1}\zeta^{-(j-1)\mathrm{p}(\mathfrak{v})}.
\end{equation*}
\end{prop}

\begin{thm}[Finite generation property]\label{thm:vertex_edge_leg_Rings}
Let $\mathrm{Cont}_{\Gamma}^{\mathrm{A}}(\mathfrak{v})$, $\mathrm{Cont}_{\Gamma}^{\mathrm{A}}(\mathfrak{e})$, and $\mathrm{Cont}_{\Gamma}^{\mathrm{A}}(\mathfrak{l})$ be as in Proposition \ref{prop:contributions}. We have $\mathrm{Cont}_{\Gamma}^{\mathrm{A}}(\mathfrak{v})\in\mathbb{C}[L^{\pm 1}]$, $\mathrm{Cont}_{\Gamma}^{\mathrm{A}}(\mathfrak{e})\in\mathbb{F}$, and $\mathrm{Cont}_{\Gamma}^{\mathrm{A}}(\mathfrak{l})\in\mathbb{F}[C_1,C_2,C_3]$. Hence, we have
\begin{equation*}
\mathcal{F}_{g, n}^{\left[\mathbb{C}^{5} / \mathbb{Z}_{5}\right]}\left(\phi_{c_{1}}, \ldots, \phi_{c_{n}}\right)\in\mathbb{F}[C_1,C_2,C_3]
\end{equation*}
and when there is no insertions, we have
\begin{equation*}
\mathcal{F}_{g}^{\left[\mathbb{C}^{5} / \mathbb{Z}_{5}\right]}\in\mathbb{F}
\end{equation*}
where $\mathbb{F}=\mathbb{C}[L^{\pm 1}][A_1,DA_1,D^2A_1,A_2]$ as before.
\end{thm}

\begin{lem}\label{lem:A2_derivative_of_edge}
We have
\begin{equation*} 
    \frac{\partial}{\partial A_{2}}\mathrm{Cont}_{\Gamma}^{\mathrm{A}}(\mathfrak{e})
    =\frac{(-1)^{b_{\mathfrak{e}1}+b_{\mathfrak{e}2}}}{5} \frac{\widetilde{P}_{3,\mathrm{p}(\mathfrak{v}_1)}^{b_{\mathfrak{e}1}}\widetilde{P}_{3,\mathrm{p}(\mathfrak{v}_2)}^{b_{\mathfrak{e}2}}}{\zeta^{(b_{\mathfrak{e}1}+3)\mathrm{p}(\mathfrak{v}_1)}\zeta^{(b_{\mathfrak{e}2}+3)\mathrm{p}(\mathfrak{v}_2)}}.
\end{equation*}
\end{lem}

\begin{thm}[The first holomorphic anomaly equation]\label{thm:HAE_wrt_A2_partial}
For $g\geq{2}$, we have
\begin{equation*}
\frac{C_{3}}{5L}\frac{\partial\mathcal{F}_{g}^{\left[\mathbb{C}^{5} / \mathbb{Z}_{5}\right]}}{\partial A_2}
=\frac{1}{2}\mathcal{F}_{g-1,2}^{\left[\mathbb{C}^5 / \mathbb{Z}_5\right]}\left(\phi_2,\phi_2\right)+\frac{1}{2}\sum_{i=1}^{g-1}\mathcal{F}_{g-i,1}^{\left[\mathbb{C}^5 / \mathbb{Z}_5\right]}\left(\phi_2\right)\mathcal{F}_{i,1}^{\left[\mathbb{C}^5 / \mathbb{Z}_5\right]}\left(\phi_2\right)
\end{equation*}
in $\mathbb{F}[C_1,C_2,C_3]$.
\end{thm}

\begin{proof}
Let $\Gamma\in\mathrm{G}_{g,0}^{\text{Dec}}(5)$ be a decorated graph and $\tilde{\mathfrak{e}}\in\mathrm{E}_{\Gamma}$ be an edge of $\Gamma$ connecting two vertices $\mathfrak{v}_1$ and $\mathfrak{v}_2$. After deleting the edge $\tilde{\mathfrak{e}}$, we obtain a new graph. (By deleting, we mean breaking the edge $\tilde{\mathfrak{e}}$ into two legs $\mathfrak{l}_{\tilde{\mathfrak{e}}}$ and $\mathfrak{l}'_{\tilde{\mathfrak{e}}}$.) There are two possibilities for the resulting graph after deletion of edge $\mathfrak{e}$:
\begin{itemize}
    \item[(i)] If it is connected, then we obtain an element of $\mathrm{G}_{g-1,2}^{\text{Dec}}(5)$, which we denote as $\Gamma_{\tilde{\mathfrak{e}}}^0$.
    
    \item[(ii)] If it is disconnected, then the resulting graph has two connected components, which we  denote as $\Gamma_{\tilde{\mathfrak{e}}}^1\in\mathrm{G}_{g_1,1}^{\text{Dec}}(5)$ and $\Gamma_{\tilde{\mathfrak{e}}}^2\in\mathrm{G}_{g_2,1}^{\text{Dec}}(5)$ where we have $g=g_1+g_2$.
\end{itemize}

By Proposition \ref{prop:contributions} and Lemma \ref{lem:A2_derivative_of_edge}, we observe that
\begin{equation*}
\begin{aligned}
\frac{\partial \mathrm{Cont}_{\Gamma}^{\mathrm{A}}(\tilde{\mathfrak{e}}) }{\partial A_2}
=&\frac{(-1)^{b_{\tilde{\mathfrak{e}}1}+b_{\tilde{\mathfrak{e}}2}}}{5} \frac{\widetilde{P}_{3,\mathrm{p}(\mathfrak{v}_1)}^{b_{\tilde{\mathfrak{e}}1}}\widetilde{P}_{3,\mathrm{p}(\mathfrak{v}_2)}^{b_{\tilde{\mathfrak{e}}2}}}{\zeta^{(b_{\tilde{\mathfrak{e}}1}+3)\mathrm{p}(\mathfrak{v}_1)}\zeta^{(b_{\tilde{\mathfrak{e}}2}+3)\mathrm{p}(\mathfrak{v}_2)}}\\
=&5\left(\frac{L^{3}}{K_{3}}\right)^2
\begin{cases}
\mathrm{Cont}_{\Gamma_{\tilde{\mathfrak{e}}}^0}^{\mathrm{A}}(\mathfrak{l}_{\tilde{\mathfrak{e}}})\mathrm{Cont}_{\Gamma_{\tilde{\mathfrak{e}}}^0}^{\mathrm{A}}(\mathfrak{l}'_{\tilde{\mathfrak{e}}})&\text{for the case (i)},\\
\mathrm{Cont}_{\Gamma_{\tilde{\mathfrak{e}}}^1}^{\mathrm{A}}(\mathfrak{l}_{\tilde{\mathfrak{e}}})\mathrm{Cont}_{\Gamma_{\tilde{\mathfrak{e}}}^2}^{\mathrm{A}}(\mathfrak{l}'_{\tilde{\mathfrak{e}}})&\text{for the case (ii)}
\end{cases}
\end{aligned}
\end{equation*}
with $\ell(\mathfrak{l}_{\tilde{\mathfrak{e}}})=\ell(\mathfrak{l}'_{\tilde{\mathfrak{e}}})=2$, i.e., with insertions $\phi_2$
.

By definition of $K_3$ and equation (\ref{eq:C12C22C3L5}), we also note that 
\begin{equation*}
\left(\frac{L^{3}}{K_{3}}\right)^2=\frac{L}{C_{3}}.
\end{equation*}
Then, for case (i), we easily see that we have
\begin{equation}\label{eq:2_insertions_graph_cont}
\begin{aligned}
\operatorname{Cont}_{\Gamma_{\tilde{\mathfrak{e}}}^0}\left(\phi_2,\phi_2\right)
=&\frac{1}{|\mathrm{Aut}(\Gamma_{\tilde{\mathfrak{e}}}^{0,\mathrm{St}})|} \sum_{\mathrm{A} \in \mathbb{Z}_{\geq 0}^{\mathrm{F}({\Gamma_{\tilde{\mathfrak{e}}}^0})}} \prod_{\mathfrak{v} \in \mathrm{V}_{\Gamma_{\tilde{\mathfrak{e}}}^0}} \mathrm{Cont}_{\Gamma_{\tilde{\mathfrak{e}}}^0}^{\mathrm{A}}(\mathfrak{v}) \prod_{\mathfrak{e} \in \mathrm{E}_{\Gamma_{\tilde{\mathfrak{e}}}^0}} \mathrm{Cont}_{\Gamma_{\tilde{\mathfrak{e}}}^0}^{\mathrm{A}}(\mathfrak{e})\prod_{\mathfrak{l} \in \mathrm{L}_{\Gamma_{\tilde{\mathfrak{e}}}^0}} \mathrm{Cont}_{\Gamma_{\tilde{\mathfrak{e}}}^0}^{\mathrm{A}}(\mathfrak{l})\\
=&\frac{1}{|\mathrm{Aut}(\Gamma_{\tilde{\mathfrak{e}}}^{0,\mathrm{St}})|} \sum_{\mathrm{A} \in \mathbb{Z}_{\geq 0}^{\mathrm{F}(\Gamma)}}\frac{C_3}{5L}\frac{\partial \mathrm{Cont}_{\Gamma}^{\mathrm{A}}(\tilde{\mathfrak{e}}) }{\partial A_2} \prod_{\mathfrak{v} \in \mathrm{V}_{\Gamma}} \mathrm{Cont}_{\Gamma}^{\mathrm{A}}(\mathfrak{v}) \prod_{\substack{
\mathfrak{e} \in \mathrm{E}_{\Gamma} \\ \mathfrak{e}\neq \tilde{\mathfrak{e}}}} \mathrm{Cont}_{\Gamma}^{\mathrm{A}}(\mathfrak{e}).
\end{aligned}
\end{equation}
Similarly, for case (ii), we observe the following
\begin{equation}\label{eq:prod_of_1_insertions_graph_conts}
\begin{aligned}
\operatorname{Cont}_{\Gamma_{\tilde{\mathfrak{e}}}^1}\left(\phi_{2}\right)&\operatorname{Cont}_{\Gamma_{\tilde{\mathfrak{e}}}^2}\left(\phi_{2}\right)\\
=&\frac{1}{|\mathrm{Aut}(\Gamma_{\tilde{\mathfrak{e}}}^{1,\mathrm{St}})|} \sum_{\mathrm{A} \in \mathbb{Z}_{\geq 0}^{\mathrm{F}({\Gamma_{\tilde{\mathfrak{e}}}^1})}} \mathrm{Cont}_{\Gamma_{\tilde{\mathfrak{e}}}^1}^{\mathrm{A}}(\mathfrak{l}_{\tilde{\mathfrak{e}}})\prod_{\mathfrak{v} \in \mathrm{V}_{\Gamma_{\tilde{\mathfrak{e}}}^1}} \mathrm{Cont}_{\Gamma_{\tilde{\mathfrak{e}}}^1}^{\mathrm{A}}(\mathfrak{v}) \prod_{\mathfrak{e} \in \mathrm{E}_{\Gamma_{\tilde{\mathfrak{e}}}^1}} \mathrm{Cont}_{\Gamma_{\tilde{\mathfrak{e}}}^1}^{\mathrm{A}}(\mathfrak{e})\\
\times&\frac{1}{|\mathrm{Aut}(\Gamma_{\tilde{\mathfrak{e}}}^{2,\mathrm{St}})|} \sum_{\mathrm{A} \in \mathbb{Z}_{\geq 0}^{\mathrm{F}({\Gamma_{\tilde{\mathfrak{e}}}^2})}} \mathrm{Cont}_{\Gamma_{\tilde{\mathfrak{e}}}^2}^{\mathrm{A}}(\mathfrak{l}'_{\tilde{\mathfrak{e}}}) \prod_{v \in \mathrm{V}_{\Gamma_{\tilde{\mathfrak{e}}}^2}} \mathrm{Cont}_{\Gamma_{\tilde{\mathfrak{e}}}^2}^{\mathrm{A}}(\mathfrak{v}) \prod_{\mathfrak{e} \in \mathrm{E}_{\Gamma_{\tilde{\mathfrak{e}}}^2}} \mathrm{Cont}_{\Gamma_{\tilde{\mathfrak{e}}}^2}^{\mathrm{A}}(\mathfrak{e})\\
=&\frac{1}{|\mathrm{Aut}(\Gamma_{\tilde{\mathfrak{e}}}^{1,\mathrm{St}})||\mathrm{Aut}(\Gamma_{\tilde{\mathfrak{e}}}^{,\mathrm{St}})|} \sum_{\mathrm{A} \in \mathbb{Z}_{\geq 0}^{\mathrm{F}(\Gamma)}}\frac{C_{3}}{5L}\frac{\partial \mathrm{Cont}_{\Gamma}^{\mathrm{A}}(\tilde{\mathfrak{e}}) }{\partial A_2} \prod_{\mathfrak{v} \in \mathrm{V}_{\Gamma}} \mathrm{Cont}_{\Gamma}^{\mathrm{A}}(\mathfrak{v}) \prod_{\substack{
\mathfrak{e} \in \mathrm{E}_{\Gamma} \\ \mathfrak{e}\neq \tilde{\mathfrak{e}}}} \mathrm{Cont}_{\Gamma}^{\mathrm{A}}(\mathfrak{e}).
\end{aligned}
\end{equation}
By Lemma \ref{lem:P0jk_is_in_CLplusminus} and Theorem \ref{thm:vertex_edge_leg_Rings}, we have the following vanishing result:
\begin{equation*}
\frac{\partial \mathrm{Cont}_{\Gamma}^{\mathrm{A}}(\mathfrak{v}) }{\partial A_2}=0
\end{equation*}
for any vertex $\mathfrak{v}\in\mathrm{V}_{\Gamma} $.

Then, this vanishing result gives us the following
\begin{equation*}
\begin{aligned}
\frac{\partial \mathrm{Cont}_{\Gamma}}{\partial A_2}
=&\frac{1}{|\mathrm{Aut}(\Gamma^{\mathrm{St}})|} \sum_{\mathrm{A} \in \mathbb{Z}_{\geq 0}^{\mathrm{F}(\Gamma)}} \prod_{\mathfrak{v} \in \mathrm{V}_{\Gamma}} \mathrm{Cont}_{\Gamma}^{\mathrm{A}}(\mathfrak{v})\frac{\partial}{\partial A_2}\left(\prod_{\mathfrak{e}\in \mathrm{E}_{\Gamma}} \mathrm{Cont}_{\Gamma}^{\mathrm{A}}(\mathfrak{e})\right)\\
=&\sum_{\tilde{\mathfrak{e}}\in \mathrm{E}_{\Gamma}}\frac{1}{|\mathrm{Aut}(\Gamma^{\mathrm{St}})|} \sum_{\mathrm{A} \in \mathbb{Z}_{\geq 0}^{\mathrm{F}(\Gamma)}}\frac{\partial \mathrm{Cont}_{\Gamma}^{\mathrm{A}}(\tilde{\mathfrak{e}}) }{\partial A_2} \prod_{\mathfrak{v} \in \mathrm{V}_{\Gamma}} \mathrm{Cont}_{\Gamma}^{\mathrm{A}}(\mathfrak{v})\prod_{\substack{
\mathfrak{e} \in \mathrm{E}_{\Gamma} \\ \mathfrak{e}\neq \tilde{\mathfrak{e}}}} \mathrm{Cont}_{\Gamma}^{\mathrm{A}}(\mathfrak{e}).
\end{aligned}
\end{equation*}
So, we have
\begin{equation}\label{eq:A2_Derivative_of_Gamma_cont}
\frac{C_{3}}{5L}\frac{\partial \mathrm{Cont}_{\Gamma}}{\partial A_2}
=\sum_{\tilde{\mathfrak{e}}\in \mathrm{E}_{\Gamma}}\frac{1}{|\mathrm{Aut}(\Gamma^{\mathrm{St}})|} \sum_{\mathrm{A} \in \mathbb{Z}_{\geq 0}^{\mathrm{F}(\Gamma)}}\frac{C_{3}}{5L}\frac{\partial \mathrm{Cont}_{\Gamma}^{\mathrm{A}}(\tilde{\mathfrak{e}}) }{\partial A_2} \prod_{\mathfrak{v} \in \mathrm{V}_{\Gamma}} \mathrm{Cont}_{\Gamma}^{\mathrm{A}}(\mathfrak{v})\prod_{\substack{
\mathfrak{e} \in \mathrm{E}_{\Gamma} \\ \mathfrak{e}\neq \tilde{\mathfrak{e}}}} \mathrm{Cont}_{\Gamma}^{\mathrm{A}}(\mathfrak{e}).
\end{equation}
Then, summing equation (\ref{eq:2_insertions_graph_cont}) and equation (\ref{eq:prod_of_1_insertions_graph_conts}) over all decorated stable graphs $\Gamma_{\tilde{\mathfrak{e}}}^0$ and $(\Gamma_{\tilde{\mathfrak{e}}}^1,\Gamma_{\tilde{\mathfrak{e}}}^2)$
we obtain
\begin{equation*}
\left\langle\left\langle\phi_{2},\phi_{2}\right\rangle\right\rangle_{g-1,2}^{\left[\mathbb{C}^5 / \mathbb{Z}_5\right]}\quad\text{and}\quad\sum_{i=1}^{g-1}\left\langle\left\langle\phi_{2}\right\rangle\right\rangle_{g-i,1}^{\left[\mathbb{C}^5 / \mathbb{Z}_5\right]}\left\langle\left\langle\phi_{2}\right\rangle\right\rangle_{i,1}^{\left[\mathbb{C}^5 / \mathbb{Z}_5\right]}
\end{equation*}
respectively. Then, by equation (\ref{eq:A2_Derivative_of_Gamma_cont}) and an analysis of automorphism factors appearing in equations (\ref{eq:2_insertions_graph_cont}), (\ref{eq:prod_of_1_insertions_graph_conts}), we conclude that
\begin{equation*}
2\frac{C_{3}}{5L}\frac{\partial}{\partial A_2}\left\langle\left\langle\right\rangle\right\rangle_{g}^{\left[\mathbb{C}^5 / \mathbb{Z}_5\right]}
=\left\langle\left\langle\phi_{2},\phi_{2}\right\rangle\right\rangle_{g-1,2}^{\left[\mathbb{C}^5 / \mathbb{Z}_5\right]}+\sum_{i=1}^{g-1}\left\langle\left\langle\phi_{2}\right\rangle\right\rangle_{g-i,1}^{\left[\mathbb{C}^5 / \mathbb{Z}_5\right]}\left\langle\left\langle\phi_{2}\right\rangle\right\rangle_{i,1}^{\left[\mathbb{C}^5 / \mathbb{Z}_5\right]}
\end{equation*}
after summing over all decorated stable graphs $\Gamma$. The reason we have $2$ in front of the left hand side is due to not having a canonical order of labelings of each of the legs $\mathfrak{l}_{\tilde{\mathfrak{e}}}$ and $\mathfrak{l}'_{\tilde{\mathfrak{e}}}$ for case (i) and double counting  of different genera of connected components for case (ii). This completes the proof. 
\end{proof}

\end{document}
